\section{Experiments}
\subsection{Training Details}
\begin{table*}[t]
  \centering
  \small
  \setlength{\tabcolsep}{8pt}
  \begin{tabular}{lll}
  \toprule
  \textbf{Parameter} & \textbf{SFT} & \textbf{GRPO} \\
  \midrule
  Training objective & Supervised fine-tuning & Group Relative Policy Optimization \\
  LoRA rank & 16 & 16 \\
  Batch size & 32 & 32 \\
  Learning rate & $1\times10^{-4}$ & $3\times10^{-5}$ \\
  Training length & 1 epoch & 50 steps \\
  Sequence / generation length & max length $=4096$ & max generation tokens $=2200$ \\
  Optimizer $\beta_1$ & 0.9 & / \\
  Optimizer $\beta_2$ & 0.95 & / \\
  Optimizer $\epsilon$ & $1\times10^{-8}$ & / \\
  Random seed & 42 & shuffle seed $=0$ \\
  Group size & / & 8 \\
  Sampling temperature & / & 1.0 \\
  Sampling concurrency & / & 8 \\
  Judge model & / & \texttt{openai/gpt-5.1} \\
  Judge concurrency & / & 32 \\
  Reward weight ($L1$) & / & 0.2 \\
  Reward weight ($L2$) & / & 0.4 \\
  Reward weight ($L3$) & / & 0.4 \\
  \bottomrule
  \end{tabular}
  \caption{Main model-agnostic hyperparameters used in SFT and GRPO. We omit dataset paths, model names, logging options, preview-saving settings,
  prompt-mode switches, and other engineering details that do not materially affect the training rule itself.}
  \label{tab:appendix_training_hparams}
  \end{table*}

Our training pipeline contains two stages: supervised fine-tuning (SFT) followed by GRPO.
  For clarity and reproducibility, we report only the model-agnostic hyperparameters that directly affect optimization, sampling, or reward
  composition, while omitting engineering-specific options such as logging, checkpoint naming, preview dumping, prompt-format switches, and task
  filtering flags.
  For SFT, we used LoRA adaptation with standard Adam-style optimization.
  For GRPO, we optimized grouped rollouts with judge-based rewards, where the final reward combined structural quality and higher-level judge scores
  using fixed weights.
  Table~\ref{tab:appendix_training_hparams} summarizes the main settings.